\section{Supported IR subset and refusal taxonomy}

Table \ref{tab:ir-subset} summarizes the constructs currently represented by
\lexir{}. ``Bounded'' means that the implementation checks a finite, declared
subset; it is not a claim that all natural-language instances of the construct
are decidable.

\begin{table}[h]
\centering
\caption{The executable boundary is explicit.}
\label{tab:ir-subset}
\small
\begin{tabularx}{\linewidth}{@{}p{0.24\linewidth}p{0.27\linewidth}X@{}}
\toprule
Construct & Status & Typical refusal or validation condition \\
\midrule
Scalar comparisons and null checks & Supported & Type mismatch or missing operand \\
Dates, durations, lists, ranges & Bounded & Unsupported temporal operator or unbounded value \\
Decision tables and hit policies & Supported with checks & Overlap, gap, or ambiguous hit policy \\
References and dependencies & Supported with checks & Unresolved or cyclic reference \\
Quantification and open-world context & Refused & \texttt{UNSUPPORTED\_SEMANTICS} with source span \\
Unbounded recursion or external side effects & Refused & \texttt{NON\_TERMINATING} or \texttt{EXTERNAL\_CONTEXT} \\
\bottomrule
\end{tabularx}
\end{table}

Refusal records contain a stable code, human-readable explanation, source
anchor, IR/schema version, and suggested review action. Validation errors are
distinct from refusals: an error means a supported construct is malformed,
whereas a refusal means the construct is outside the declared language.
